\chapter{Terminology}

\begin{reader-guide}
 Feel free to skip this chapter if you are a first-time reader. 
\end{reader-guide}

\index{Requirements!Non-functional}
Performance and efficiency, i.e.~the speed of a code, are critical non-functional properties of scientific codes.
Measureing, asessing and interpreting them is challenging; not even starting to
think about tuning.
Most annoyingly, there is no standardised way how to analyse
and assess performance.
There is no vanilla workflow.
Every team, every project, every scientist create their own analysis workflow or
process, and this processes often changes from paper to paper or experiment to experiment.

Despite the vagueness in the metrics and the processes, there are some
well-defined, common terms in the community that we should use.
There are also some key concepts and some jargon that we use throughout the
write-up.
So we better introduce them first.


